% Original analytic teaching diagram, 2026-11-09.
\definecolor{VB}{HTML}{7188A3}\definecolor{VBF}{HTML}{E5EAF0}
\definecolor{VG}{HTML}{829C8B}\definecolor{VGF}{HTML}{E8EEE9}
\definecolor{VY}{HTML}{B4A16F}\definecolor{VYF}{HTML}{F1ECDD}
\definecolor{VR}{HTML}{A7787D}\definecolor{VRF}{HTML}{F0E5E6}
\definecolor{VP}{HTML}{9588A5}\definecolor{VPF}{HTML}{ECE8F0}
\definecolor{VI}{HTML}{3B4252}\definecolor{VM}{HTML}{6B7280}\definecolor{VL}{HTML}{4C4D4F}
\begin{tikzpicture}[x=1mm,y=1mm,>=Stealth,font=\figurefont\fontsize{8}{10}\selectfont,
lab/.style={text=VM,font=\figurefont\fontsize{7.5}{9}\selectfont},
box/.style={rounded corners=1.2mm,draw=VB,fill=white,line width=.8pt,minimum height=10mm,inner sep=1.2mm,align=center,text=VI},
chk/.style={rounded corners=1.2mm,draw=VP,fill=VPF,line width=.8pt,minimum height=10mm,inner sep=1.2mm,align=center,text=VI},
arr/.style={-{Stealth[length=2mm]},VL,line width=.8pt}]
\node[box,minimum width=17mm] (p) at (0,14) {原问题};
\node[box,minimum width=22mm,fill=VGF] (e) at (29,26) {等价改写};
\node[box,minimum width=22mm,fill=VYF] (r) at (29,2) {半正定松弛};
\node[box,minimum width=25mm] (s) at (60,26) {求解\\丢弃辅助变量};
\node[box,minimum width=25mm] (t) at (60,2) {求解松弛\\恢复与取整};
\node[chk,minimum width=20mm] (c) at (91,14) {分别核验};
\draw[arr] (p.east)--(e.west); \draw[arr] (p.east)--(r.west);
\draw[arr] (e)--(s); \draw[arr] (r)--(t);
\draw[arr] (s.east)--(c.west); \draw[arr] (t.east)--(c.west);
\node[lab,anchor=north] at (91,8) {可行解给上界};
\node[lab,anchor=north] at (91,4) {对偶可行给下界};
\end{tikzpicture}
